\section{Concluding remarks}

The results reported in this paper shed some new light
on the nature of non-abelian gauge theories
and the continuum limit in lattice QCD.
They are obviously incomplete in several respects
and give rise to some interesting questions.
In particular, an all-order analysis of the Wilson flow in
perturbation theory will probably be required in order to achieve a
structural understanding of its renormalization properties.

The one-loop calculation in sect.~2 suggests that
the fields obtained at positive flow time are renormalized
fields for any number of quark flavours.
So far the question was studied numerically
only in the pure gauge theory,
but it may be encouraging to note that the conjecture is true
in QED, for any charged matter multiplet,
as a consequence of the gauge Ward identity.

An intriguing aspect of the transformed functional
integral~\eqref{eq:fint} is the fact that it is dominated by
smooth fields. Since the field transformation~\eqref{eq:fieldtrafo} is invertible,
the integral nevertheless encodes all the physics
described by the theory.
Some properties
(the division into topological sectors,
for example) however become particularly transparent in this formulation
of the theory,
while its behaviour at high energies is better
discussed in terms of the fundamental gauge field.
